% =====================================================================
%  Robotnik DAgger theory -- drop-in section (no preamble, no document env)
%  Usage:     \input{dagger_section}
%  Needs in the main preamble (nothing else):
%     \usepackage{amsmath,amssymb,booktabs,tikz,graphicx}
%  Figures are read from \dagfigdir (default: figures/ next to this file);
%  \newcommand{\dagfigdir}{path/to/figures/} before \input to change it.
%     \usetikzlibrary{positioning,arrows.meta,calc,backgrounds}
%  Change \section to \subsection / \chapter if needed.
%  Labels are prefixed "dag:" ; helper macros use \providecommand.
%  Verified against: dagger_common.py, robotnik_arm_control.py, arm_v4.py,
%  robotnik_dagger_collect.py, train_dagger.py, train_arm_v5.py, README_V53.
% =====================================================================

% ---------- local macros ----------
\providecommand{\dagE}{\mathcal{E}}                        % expert
\providecommand{\dagPi}{\pi_\theta}                        % learner policy
\providecommand{\dagD}{\mathcal{D}}                        % dataset
\providecommand{\dagnorm}[1]{\lVert #1\rVert_2}            % Euclidean norm
\providecommand{\dagwrap}{\operatorname{wrap}}             % angle wrapping
\providecommand{\dagNm}{\mathrm{N\!\cdot\!m}}              % newton-metre
\providecommand{\dagone}{\mathbf{1}}                       % indicator
\providecommand{\dagfigdir}{figures/}                      % figure folder

% ---------- colours and a light-weight "box" (no extra packages) ----------
\definecolor{dagblue}{HTML}{1F4E79}
\definecolor{dagblueL}{HTML}{E8F0F8}
\definecolor{daggreen}{HTML}{2E7D32}
\definecolor{daggreenL}{HTML}{E6F4E8}
\definecolor{dagorange}{HTML}{C26A00}
\definecolor{dagorangeL}{HTML}{FFF1DC}
\providecommand{\dagbox}[3]{% {frame colour}{fill colour}{content}
  \par\medskip\noindent
  \fcolorbox{#1}{#2}{\begin{minipage}{\dimexpr\linewidth-2\fboxsep-2\fboxrule\relax}%
    \small\setlength{\parskip}{0.3em}#3\end{minipage}}\par\medskip}

% =====================================================================
\section{DAgger Imitation Learning for the Robotnik Mobile Manipulator}
\label{dag:sec}

% ---------------------------------------------------------------------
\subsection{Problem and key idea}
\label{dag:problem}

A simulated Robotnik mobile manipulator (PyBullet physics, Pinocchio kinematics and dynamics) must (i)~drive its four-wheel base to a standoff pose in front of a world-frame target and then (ii)~move its 7-joint arm so that the end-effector reaches the target and stays there. At step $t$ the robot is in state $s_t$ and receives a command $a_t=[u_{wheel},u_{arm}]$, where $u_{wheel}\in\mathbb{R}^4$ are wheel angular velocities (rad/s) and $u_{arm}\in\mathbb{R}^7$ are joint torques ($\dagNm$). Learned policies $\dagPi$ with trainable weights $\theta$ are trained to imitate a model-based \emph{expert controller} $\dagE$:
\begin{equation}
  \dagPi\bigl(\phi(s_t)\bigr)\;\approx\;a_t^{*}=\dagE(s_t),
  \label{dag:goal}
\end{equation}
where $\phi(s_t)$ is a fixed-length feature vector (Sec.~\ref{dag:state}); we write $\dagPi(s_t)$ for $\dagPi(\phi(s_t))$.

Training only on states the expert visits (\emph{behaviour cloning}) is not enough: small learner errors lead to states the expert never visited, where the network errs more and the errors \emph{compound}. In the classical analysis (Ross et al., 2011), for a horizon of $T$ steps and per-step imitation error $\epsilon$, behaviour cloning can accumulate $O(T^{2}\epsilon)$ error while DAgger achieves $O(T\epsilon)$. This guarantee applies only to a policy that is itself trained with DAgger on the distribution it induces; in this project that is the arm (see below), whereas the wheel network is pure behaviour cloning.

\dagbox{dagblue}{dagblueL}{%
\textbf{Key idea of DAgger (Dataset Aggregation).} Let the \emph{learner} drive, ask the \emph{expert} what it would have done in every state the learner actually reaches, add these expert-labelled states to the training set, and retrain.}

\textbf{What is DAgger in this project.} The base is driven by a 38-input, raw-feature wheel network trained once by behaviour cloning (\texttt{train\_dagger.py}) on the 1400 training episodes of the 2000 original expert demonstrations, and copied unchanged into every later checkpoint; it is never retrained by DAgger. The DAgger rounds collect data for the \emph{arm} networks, which are retrained from scratch on the original demonstrations plus all aggregated rollouts (Secs.~\ref{dag:collect} and~\ref{dag:train}).

% ---------------------------------------------------------------------
\subsection{Notation}
\label{dag:notation}

\begin{center}
\small
\renewcommand{\arraystretch}{1.12}
\begin{tabular}{@{}p{3.3cm}p{9.4cm}@{}}
\toprule
\textbf{Symbol} & \textbf{Meaning}\\
\midrule
$s_t,\ \phi(s_t)$ & robot state and its 38-dimensional feature vector at step $t$\\
$\phi_{arm}(t)$ & 66-dimensional arm input: $\phi(s_t)$ plus joint history at lags 10 and 30 steps\\
$a_t,\ a_t^{*}$ & executed and expert (target) action, $\in\mathbb{R}^4\times\mathbb{R}^7$\\
$\dagE,\ \dagPi$ & expert controller; learner policy with weights $\theta$\\
$a_t^{expert},a_t^{learner}$ & actions proposed by $\dagE$ and $\dagPi$ on the same state\\
$\beta,\ M$ & probability of expert control per mixing interval; interval length in steps\\
$\dagD,\ N$ & aggregated dataset of pairs $(\text{input},a^{*})$ and its size\\
$d_{\dagPi}$ & state distribution induced by running $\dagPi$\\
$p_{base},\theta_{base}$ & base position and yaw\\
$p_{goal},\theta_{goal}$ & base goal position and goal bearing (derived from the target)\\
$p_{ee},\ p_{target}$ & end-effector position and world-frame target\\
$d,\ \Delta\theta_{head},\Delta\theta_{ori}$ & distance to goal, heading error toward the goal, final orientation error\\
$c_t,\ \hat c_t$ & raw and latched navigation-complete flags\\
$q,\dot q,\ q_{init}$ & arm joint angles, velocities (7 joints) and initial arm configuration\\
$q_{des},\dot q_{des}$ & desired joint angles and velocities (expert reference)\\
$J,\ v_{cart},\ \lambda$ & end-effector Jacobian ($3\times7$), desired Cartesian velocity, IK damping\\
$K_p,K_d,\ g(q)$ & PD gains and gravity-compensation torque\\
$\varepsilon,\ N_{hold}$ & reach tolerance and required number of consecutive hold steps\\
\bottomrule
\end{tabular}
\end{center}

% ---------------------------------------------------------------------
\subsection{Task geometry and state features}
\label{dag:state}

\paragraph{Base goal.} The target $p_{target}\in\mathbb{R}^3$ is given in the world frame. The base is not sent to the target itself but to a \emph{standoff} pose on the ray from the world origin through the target's $xy$-position:
\begin{equation}
  \theta_{goal}=\operatorname{atan2}(p_{target,y},p_{target,x}),\qquad
  p_{goal}=p_{target,xy}-\sigma(p_{target,z})\begin{bmatrix}\cos\theta_{goal}\\ \sin\theta_{goal}\end{bmatrix},
  \label{dag:goalpose}
\end{equation}
where the standoff distance $\sigma(z)$ is piecewise-linear in the target height, interpolating $(z,\sigma)=(0.1,0.80),(0.5,0.65),(1.1,0.65),(1.25,0.35)$ (m) and constant outside this range. Higher targets are approached more closely.

\paragraph{Navigation signals.} From the base pose the code computes
\begin{equation}
\begin{aligned}
  d&=\dagnorm{p_{goal}-p_{base}},\qquad
  \Delta\theta_{head}=\dagwrap\!\bigl(\operatorname{atan2}(\Delta p_y,\Delta p_x)-\theta_{base}\bigr),\\
  \Delta\theta_{ori}&=\dagwrap(\theta_{goal}-\theta_{base}),\qquad \Delta p=p_{goal}-p_{base},
\end{aligned}
\label{dag:geom}
\end{equation}
where $\dagwrap(\cdot)$ maps an angle to $(-\pi,\pi]$ so that the shortest rotation is used ($179^\circ$ versus $-179^\circ$ is a $2^\circ$ error, not $358^\circ$). $\Delta\theta_{head}$ says how far the base must turn to \emph{face the goal point}; $\Delta\theta_{ori}$ says how far it must turn to reach the final goal \emph{orientation}. The observable \emph{navigation-complete} flag is
\begin{equation}
  c_t=\dagone\bigl[(d<0.08\ \text{m}\ \wedge\ |\Delta\theta_{ori}|<0.05\ \text{rad})\ \vee\ t_{elapsed}>25\ \text{s}\bigr],
  \label{dag:nav}
\end{equation}
so a 25\,s timeout also forces completion. In practice this timeout term is never active in the data: the demonstrations last 8\,s and the collector stops after 20\,000 steps (20\,s at $\Delta t=1$\,ms), so in all training rows and in the collector-based validation runs $c_t$ reduces to the position-and-orientation test. Only the 30\,s deployment runner can reach it.

\paragraph{Feature vector.} The 38 raw (unscaled) features are: base position $(x,y,z)$ and yaw (4); target position (3); end-effector position (3); goal position $(x,y)$ and goal bearing (3); arm joint angles $q$ (7) and velocities $\dot q$ (7); initial arm angles $q_{init}$ (7); and the four navigation signals $d,\Delta\theta_{head},\Delta\theta_{ori},c_t$ (4). A metadata file stores the feature and output names of every checkpoint and the loader refuses a mismatch, so training and deployment inputs always have identical meaning. Only quantities observable at deployment are provided: the expert's internal references (e.g.\ $q_{des}$) are logged for diagnosis but \emph{never} given to the learner.

% ---------------------------------------------------------------------
\subsection{Expert controller}
\label{dag:expert}

The expert is a \emph{stateful} controller: the base keeps a phase variable and the arm keeps its reference $q_{des}$. Hence $a_t^{expert}=\dagE(s_t;z_t)$ with internal state $z_t=(\text{base phase},q_{des})$. To keep $z_t$ consistent, the collector advances the expert on \emph{every} observed state, including steps where the learner is the one acting.

\paragraph{Base (skid-steer).} The phases are \emph{align}~$\to$~\emph{drive}~$\to$~\emph{orient}~$\to$~\emph{done}:
\begin{itemize}
  \item \emph{align}: turn in place with $\omega=\operatorname{clip}(3\,\Delta\theta_{head},\pm1.0)$ rad/s until $|\Delta\theta_{head}|<0.05$ (or $d<0.08$);
  \item \emph{drive}: speed $v=\operatorname{clip}(1.0\,d,\,0.05,\,0.25)\cdot\max(0,\cos\Delta\theta_{head})$ m/s and gentle steering $\omega=\operatorname{clip}(1.5\,\Delta\theta_{head},\pm0.3)$ rad/s, until $d<0.08$ or the goal has been passed ($d<0.15$ and $\cos\Delta\theta_{head}<0$);
  \item \emph{orient}: turn with $\omega=\operatorname{clip}(3\,\Delta\theta_{ori},\pm1.0)$ until $|\Delta\theta_{ori}|<0.05$, then \emph{done}: wheels are held at zero. A 25\,s timeout also sets \emph{done}.
\end{itemize}
Forward speed $v$ and yaw rate $\omega$ become wheel speeds $\omega_{L,R}=(v\mp\omega B/2)/r$ with wheel radius $r=0.11$\,m and effective track width $B=0.66$\,m; the four commands are ordered $[\omega_R,\omega_L,\omega_L,\omega_R]$ and clipped to $\pm3.18$ rad/s at execution.

\paragraph{Arm: two modes.} While the base is not done, the arm \emph{holds} its initial configuration: $q_{des}=q_{init}$, $\dot q_{des}=0$. Once the base is done, the arm \emph{reaches}. The reach mode first clamps the target to a sphere of radius $0.85$\,m around the shoulder (link~2), $p'_{target}$, and commands a proportional Cartesian velocity capped in norm:
\begin{equation}
  v_{cart}=\operatorname{clip}_{\lVert\cdot\rVert\le0.25\,\text{m/s}}\bigl(1.5\,(p'_{target}-p_{ee})\bigr).
  \label{dag:vcart}
\end{equation}
The Jacobian $J\in\mathbb{R}^{3\times7}$ satisfies $v_{ee}=J\dot q$. Because $J$ is not square and loses rank near singular poses, the expert solves a \emph{damped least-squares} problem that trades tracking accuracy against joint speed:
\begin{equation}
  \dot q_{des}=\arg\min_{\dot q}\ \dagnorm{J\dot q-v_{cart}}^{2}+\lambda^{2}\dagnorm{\dot q}^{2}
  =J^{\top}\bigl(JJ^{\top}+\lambda^{2}I_3\bigr)^{-1}v_{cart}.
  \label{dag:dls}
\end{equation}
The closed form follows from setting the gradient to zero, $(J^{\top}J+\lambda^{2}I)\dot q=J^{\top}v_{cart}$, and the identity $(J^{\top}J+\lambda^{2}I)^{-1}J^{\top}=J^{\top}(JJ^{\top}+\lambda^{2}I)^{-1}$. The damping adapts to the smallest singular value $\sigma_{\min}(J)$: $\lambda$ is interpolated linearly from $0.10$ (at $\sigma_{\min}=0$) to $0.01$ (at $\sigma_{\min}\ge0.15$), so it is large near singularities. Joints within $0.05$\,rad of a limit that the solution would push further into it are frozen, and (\ref{dag:dls}) is re-solved over the remaining joints (up to 7 passes). The result is scaled so that $\max_j|\dot q_{des,j}|\le1.0$ rad/s, and the reference is \emph{integrated and clipped}:
\begin{equation}
  q_{des}\leftarrow q_{des}+\dot q_{des}\,\Delta t,\qquad
  q_{des}\in[\,q-q_{lead},\,q+q_{lead}\,],\quad q_{lead}=0.15\ \text{rad},
  \label{dag:qlead}
\end{equation}
and finally kept $0.05$\,rad inside the joint limits. The window $q_{lead}$ is an anti-windup measure: for an unreachable target the reference cannot run away from the measured arm.

\paragraph{Torque law.} In both modes a PD law with gravity compensation produces the torques, clipped \emph{after} adding gravity:
\begin{equation}
  \tau_{arm}=\operatorname{clip}_{[-30,30]\,\dagNm}\Bigl(K_p(q_{des}-q)+K_d(\dot q_{des}-\dot q)+g(q)\Bigr),\quad K_p=100,\ K_d=30.
  \label{dag:pd}
\end{equation}
The $K_p$ term is a spring toward $q_{des}$, the $K_d$ term a damper, and $g(q)$ (from the Pinocchio dynamics model) cancels gravity. The gripper fingers are held at a fixed $0.1$\,rad by position control and are not part of the policy.

% ---------------------------------------------------------------------
\subsection{Learned policy: three networks}
\label{dag:policy}

The deployed policy consists of \emph{three} MLPs:
\begin{itemize}
  \item \textbf{wheel network}: $\phi(s_t)\in\mathbb{R}^{38}\to u_{wheel}\in\mathbb{R}^4$ (behaviour-cloned once and kept fixed, Sec.~\ref{dag:train});
  \item \textbf{arm HOLD network} and \textbf{arm REACH network}: $\phi_{arm}(t)\in\mathbb{R}^{66}\to u_{arm}\in\mathbb{R}^7$.
\end{itemize}
The arm input appends a short joint-state history so the network can infer motion trends,
\begin{equation}
  \phi_{arm}(t)=\bigl[\ \phi(s_t);\ y_{\max(0,t-10)};\ y_{\max(0,t-30)}\ \bigr],\qquad y_t=(q_t,\dot q_t)\in\mathbb{R}^{14},
  \label{dag:lags}
\end{equation}
giving $38+14+14=66$ inputs. The lags are 10 and 30 \emph{rows} of the trajectory and are clamped at the start of an episode. At deployment and in the DAgger rollouts one row is one simulator step (1\,ms), so the lags are 10\,ms and 30\,ms. The original demonstrations, however, are logged every 10\,ms (801 rows per 8\,s episode), so on $\dagD_0$ the same lags span 100\,ms and 300\,ms. The history features therefore do not have the same physical meaning in $\dagD_0$ as at deployment; this is a known limitation of the current pipeline. The arm network is selected by the \emph{latched} navigation flag,
\begin{equation}
  \hat c_t=\max_{\tau\le t}c_\tau,\qquad
  \text{use }\begin{cases}\text{REACH network}&\hat c_t=1,\\ \text{HOLD network}&\hat c_t=0,\end{cases}
  \label{dag:latch}
\end{equation}
and at deployment $\hat c_t$ replaces the raw flag in the feature vector of all three networks. The arm networks are also trained with $\hat c_t$; the wheel network, however, was trained with the raw flag $c_t$, so for it the latched input is a (small) train/deployment difference. The running maximum keeps the phase at 1 once navigation has finished, even if the raw flag later flickers back to 0. \textbf{Note on naming:} \emph{HOLD} is the phase \emph{before} navigation completes (the arm holds its home pose $q_{init}$), and \emph{REACH} is the phase after it (the arm moves to the target). This corresponds to the expert modes above, with one approximation: the expert switches to reach on its internal flag (\emph{done} of the base state machine), which, via the passed-goal shortcut, can in principle become true at $0.08\le d<0.15$\,m while $\hat c_t=0$. In a check of 40 DAgger rollouts (about 262\,000 rows) the two flags never disagreed.

Each network has the form
\begin{multline*}
  \mathrm{Lin}(d_{in},256)\to\mathrm{ReLU}\to\mathrm{Drop}(0.1)\to\mathrm{Lin}(256,256)\to\mathrm{ReLU}\to\mathrm{Drop}(0.1)\\
  \to\mathrm{Lin}(256,128)\to\mathrm{ReLU}\to\mathrm{Drop}(0.1)\to\mathrm{Lin}(128,d_{out}),
\end{multline*}
where $\mathrm{Lin}(m,n)$ is a fully connected layer $y=Wx+b$, $\mathrm{ReLU}(x)=\max(0,x)$ supplies the nonlinearity, and $\mathrm{Drop}(0.1)$ zeroes $10\%$ of activations during training only. No feature scaling is used (raw inputs).

% ---------------------------------------------------------------------
\subsection{DAgger data collection}
\label{dag:collect}

At every step $t$ of a rollout the loop in Fig.~\ref{dag:fig} runs: (1)~compute $\phi(s_t)$ and the learner action $a_t^{learner}$; (2)~advance the expert and obtain $a_t^{expert}$ on the same state; (3)~execute the action selected by the source setting; (4)~store the \emph{expert} action as the label; (5)~step the simulator.

The wheels and the arm each have a source $\in\{\text{expert},\text{learner},\text{mix}\}$. In \emph{mix} the choice is made per \emph{interval}, not per step: every $M$ steps (default $M=100$) one Bernoulli draw $b_k\sim\mathrm{Bernoulli}(\beta)$ is made, and
\begin{equation}
  a_t=\begin{cases}a_t^{expert}&\text{if }b_{\lfloor t/M\rfloor}=1,\\ a_t^{learner}&\text{otherwise,}\end{cases}
  \label{dag:mix}
\end{equation}
where $\beta$ is the probability that the expert controls the interval (CLI default $0.8$) and the same draw is shared by wheels and arm when both are in mix. Holding the choice for $M$ steps lets a learner segment run long enough to drift into realistic error states. The command modes are \emph{collect} (data generation), \emph{expert}, \emph{learner}, \emph{deploy} (learner only, the expert is never called) and \emph{diagnose} (hybrid tests).

\paragraph{Rounds actually used.} The first arm-DAgger round used expert wheels and a mixed arm ($\beta=0.5$, $M=50$). The V5.3-A round used expert wheels, learner arm and $\beta=0$ (pure learner arm), on the fixed targets of \texttt{targets.csv} with the default base spawn. V5.3-A was trained on the union of both rounds (51 + 110 rollouts, of which 160 belong to training IDs). A third round (V5.4, Sec.~\ref{dag:v54}) added rollouts with a \emph{random target and random base spawn} per episode. Rollouts are generated for the training episode IDs only; rollouts of validation/test IDs are excluded from training. After each round the dataset is aggregated,
\begin{equation}
  \dagD_{k+1}=\dagD_{k}\cup\bigl\{(\phi_{arm}(t),\,a_t^{expert})\ :\ \text{states visited in round }k\bigr\},
  \label{dag:agg}
\end{equation}
where $\dagD_0$ holds the original demonstrations (2000 episodes split by episode into 1400/300/300 train/validation/test), and the arm networks are retrained from scratch on $\dagD_{k+1}$. The label is \emph{always} the expert's, whichever policy drove the robot, so the supervision stays correct while the learner is in control. Rows are kept in temporal order inside each rollout because the lag features depend on it.

\begin{figure}[ht]
\centering
\begin{tikzpicture}[
  box/.style={draw=dagblue, rounded corners, align=center, fill=dagblueL, minimum height=0.9cm, font=\small},
  proc/.style={draw=daggreen, rounded corners, align=center, fill=daggreenL, minimum height=0.9cm, font=\small},
  data/.style={draw=dagorange, rounded corners, align=center, fill=dagorangeL, minimum height=0.9cm, font=\small},
  arr/.style={-{Latex[length=2.5mm]}, thick}]
\node[box]  (s) {State $s_t$\\$\phi(s_t)$};
\node[proc, above right=0.5cm and 1.2cm of s] (l) {Learner\\$\dagPi$};
\node[proc, below right=0.5cm and 1.2cm of s] (e) {Expert\\$\dagE(s_t;z_t)$};
\node[box, right=4.6cm of s] (x) {Execute $a_t$\\(source, $\beta$, $M$)};
\node[data, right=1.2cm of x] (d) {Dataset $\dagD$\\$\leftarrow (\text{input},a_t^{expert})$};
\draw[arr] (s.east)--(l.west);
\draw[arr] (s.east)--(e.west);
\draw[arr] (l.east)--(x.north west);
\draw[arr] (e.east)--(x.south west);
\draw[arr] (x)--(d);
\draw[arr] (e.south) -- ++(0,-0.6) -| node[pos=0.25,below,font=\footnotesize]{expert label} (d.south);
\draw[arr] (x.north) -- ++(0,1.9) -| node[pos=0.25,above,font=\footnotesize]{simulator step: next state $s_{t+1}$} (s.north);
\end{tikzpicture}
\caption{DAgger data collection. Learner and expert are evaluated on the same state (the stateful expert is advanced at every step); the source setting decides which action is executed, and the expert action is always stored as the label.}
\label{dag:fig}
\end{figure}

% ---------------------------------------------------------------------
\subsection{Objective and training}
\label{dag:train}

DAgger targets the state distribution the learner itself produces,
\begin{equation}
  \min_{\theta}\ \mathbb{E}_{s\sim d_{\dagPi}}\bigl[\ \ell\bigl(\dagPi(s),\dagE(s)\bigr)\ \bigr],
  \label{dag:obj}
\end{equation}
with $\ell$ the squared error. Since $d_{\dagPi}$ depends on $\theta$, it is approximated iteratively by the collect--aggregate--retrain loop above (the expert's internal state $z_t$ is not observed by the learner, so the match is approximate). A second approximation: the arm rollouts were collected with \emph{expert} wheels, so the aggregated data cover the distribution induced by ``expert base + learner arm'', not the deployed ``learner base + learner arm''. Deviations of the learned base trajectory are therefore not corrected by DAgger; the remaining navigation failures (Sec.~\ref{dag:task}) are consistent with this.

\paragraph{Wheel network} (behaviour cloning, 1400 training episodes, raw-flag input $c_t$): mean squared error over output components, AdamW (lr $10^{-3}$, weight decay $10^{-5}$), $\le200$ epochs with early stopping (patience 20) and learning-rate reduction on plateau; the optional smoothness weight defaults to $0$.

\paragraph{Arm networks.} For network $\mathrm{phase}\in\{\text{HOLD},\text{REACH}\}$, only rows with the matching latched flag are used (HOLD: $\hat c=0$; REACH: $\hat c=1$). With input $x_i=\phi_{arm}(t_i)$ and expert torque $\tau_i^{*}$, the loss is a weighted mean squared error plus a temporal-smoothness term:
\begin{equation}
  \mathcal{L}(\theta)=\frac1N\sum_{i=1}^{N}\Bigl[\,w_i\,\tfrac{1}{7}\dagnorm{\dagPi(x_i)-\tau_i^{*}}^{2}
  +\lambda_{smooth}\,m_i\,\tfrac{1}{7}\dagnorm{\dagPi(x_i)-\dagPi(x_{i-1})}^{2}\Bigr],
  \label{dag:loss}
\end{equation}
where $x_{i-1}$ is the previous row of the same episode (1\,ms earlier in DAgger rollouts, 10\,ms earlier in the demonstrations), $m_i=1$ only if that row exists and belongs to the same phase (else $0$), and $\lambda_{smooth}=0.01$ discourages jerky torques. The \emph{recovery weight}
\begin{equation}
  w_i=1+4\cdot\dagone\bigl[\dagnorm{p_{ee}-p_{target}}<0.35\ \text{m}\ \wedge\ \text{phase}=\text{REACH}\bigr]
  \label{dag:weight}
\end{equation}
gives samples near the target five times the weight, so the REACH network learns to stabilise at the target and not only to approach it. Training uses AdamW (lr $5\times10^{-4}$, weight decay $10^{-5}$), batch size 1024, gradient-norm clipping at 1.0, up to 100 epochs with early stopping (patience 15) on the \emph{unweighted} validation MSE, and at most 3000 randomly sampled rows per episode and phase.

% ---------------------------------------------------------------------
\subsection{Task phases, success criterion and stopping}
\label{dag:task}

The task has a natural order (Fig.~\ref{dag:phases}): the base navigates while the arm holds its home pose; once navigation completes, the arm reaches and must then stay on target.

\begin{figure}[ht]
\centering
\begin{tikzpicture}[
  ph/.style={draw, rounded corners, align=center, minimum height=1.0cm, minimum width=3.7cm, font=\footnotesize},
  arr/.style={-{Latex[length=2.5mm]}, thick}]
\node[ph, fill=dagblueL, draw=dagblue] (a) {1. Navigate (arm HOLD)\\align $\to$ drive $\to$ orient};
\node[ph, fill=daggreenL, draw=daggreen, right=1.3cm of a] (b) {2. Reach (arm REACH)\\arm moves to $p_{target}$};
\node[ph, fill=dagorangeL, draw=dagorange, right=1.9cm of b] (c) {3. Success\\$\dagnorm{p_{ee}-p_{target}}<\varepsilon$};
\draw[arr] (a)--(b) node[midway,above=2pt,font=\scriptsize]{$\hat c=1$};
\draw[arr] (b)--(c) node[midway,above=2pt,font=\scriptsize]{$N_{hold}$};
\end{tikzpicture}
\caption{Task phases and the network used in each. The latched flag $\hat c$ (Eq.~\eqref{dag:latch}) switches the arm from HOLD to REACH.}
\label{dag:phases}
\end{figure}

A rollout \emph{succeeds} if the target is held, not merely touched:
\begin{equation}
  \dagnorm{p_{ee}-p_{target}}<\varepsilon\ \ \wedge\ \ \text{navigation done}\qquad\text{for }N_{hold}\text{ consecutive steps},
  \label{dag:succ}
\end{equation}
with $\varepsilon=0.05$\,m and $N_{hold}=30$ steps (about 30\,ms at 1\,ms per step); at deployment ``navigation done'' is the position-and-orientation test of Eq.~\eqref{dag:nav} without the timeout. Other stop reasons are the step cap (20\,000 steps in the collector, 30\,000 in the runner), a joint-angle guard ($|q_j|>15$ rad) and a joint-velocity guard ($|\dot q_j|>8$ rad/s). The stall-detection class \texttt{StopCriteria} exists in the control module but is not called by the collector. Reported metrics are the final and minimum end-effector error, the fraction of steps within $\varepsilon$, time of first entry into the tolerance band, sustained-reach duration, and CPU/GPU utilisation.

\paragraph{Reported result.} On the targets of the 300 validation episodes (run from the default base spawn and arm configuration; the per-episode base and arm offsets stored in \texttt{targets.csv} are not re-applied by the collector) the V5.3-A arm reaches 280 successes (93.3\%) versus 263 (87.7\%) for the previous version; the 300 test episodes are still untouched. Remaining failures include episodes where the base gets within about 6\,cm of its goal but never satisfies both navigation tolerances, and frequent joint-6 velocity jumps.

% ---------------------------------------------------------------------
\subsection{Third DAgger round (V5.4) and evaluation}
\label{dag:v54}

\paragraph{Round 3.} The V5.3-A policy was rolled out with expert wheels and the learner arm ($\beta=0$) on 1400 episodes, each with a random target and a random base spawn. Target $x,y,z$ were drawn uniformly from the per-axis ranges of \texttt{targets.csv} ($x,y\in[-1.8,1.8]$\,m, $z\in[0.3,0.9]$\,m); the spawn offset was drawn from the demonstration offsets ($x,y\in[-0.25,0.25]$\,m, any yaw). Because the targets are random, the episode IDs of these rollouts are labels only. Aggregating with rounds 1--2 gives $\dagD_3=\dagD_0\cup$ 1561 rollouts (51 + 110 + 1400). The arm networks were retrained from scratch on $\dagD_3$ with the hyperparameters of Sec.~\ref{dag:train}, and the wheel network was copied unchanged. The resulting checkpoint is called V5.4.

\paragraph{Protocol.} Both checkpoints were evaluated closed-loop (learner wheels and learner arm, success as in Eq.~\eqref{dag:succ}) on (i)~the 300 validation targets from the default spawn with a 20\,000-step cap, and (ii)~100 episodes with a random target and random spawn drawn from the round-3 distribution, using the same five seeds for both models so that they face identical episodes. Intervals are 95\,\% Wilson score intervals.

\begin{center}
\small
\renewcommand{\arraystretch}{1.12}
\begin{tabular}{@{}lcccc@{}}
\toprule
 & \multicolumn{2}{c}{300 validation targets} & \multicolumn{2}{c}{100 random target + spawn}\\
 & V5.3-A & V5.4 & V5.3-A & V5.4\\
\midrule
Successes & \textbf{280} (93.3\,\%) & 265 (88.3\,\%) & \textbf{76} & 75\\
95\,\% interval & [89.9, 95.6]\,\% & [84.2, 91.5]\,\% & [66.8, 83.3]\,\% & [65.7, 82.5]\,\%\\
Step-cap / velocity-guard stops & 20 / 0 & 35 / 0 & 24 / 0 & 23 / 2\\
Median final EE error & 46.7\,mm & 47.1\,mm & 47.6\,mm & 47.8\,mm\\
Mean time to success & 6.78\,s & 6.70\,s & 7.22\,s & 6.95\,s\\
Best validation MSE, HOLD / REACH & 0.032 / 0.300 & 0.080 / 0.469 & -- & --\\
\bottomrule
\end{tabular}
\end{center}

\begin{figure}[ht]
\centering
\includegraphics[width=\linewidth]{\dagfigdir v54_training_curves.pdf}
\caption{Validation MSE of the arm networks (teacher-forced, on the demonstration validation split, log scale); dots mark the saved checkpoints. V5.4 converges to a higher error on both networks.}
\label{dag:fig:v54train}
\end{figure}

\begin{figure}[ht]
\centering
\includegraphics[width=0.62\linewidth]{\dagfigdir v54_success_rate.pdf}
\caption{Closed-loop success rate with 95\,\% Wilson intervals.}
\label{dag:fig:v54succ}
\end{figure}

\begin{figure}[ht]
\centering
\includegraphics[width=\linewidth]{\dagfigdir v54_final_error_cdf.pdf}
\caption{Cumulative distribution of the final end-effector error (log scale); the dashed line is the 5\,cm tolerance. Successful episodes stop just inside the tolerance (the hold criterion ends the episode), so both curves jump at about 47\,mm; the right-hand tails are step-cap failures.}
\label{dag:fig:v54cdf}
\end{figure}

\begin{figure}[ht]
\centering
\includegraphics[width=\linewidth]{\dagfigdir v54_stop_reasons.pdf}
\caption{Why episodes ended. Almost every failure hits the step cap; the joint-velocity guard fired only twice (V5.4, random spawns).}
\label{dag:fig:v54stop}
\end{figure}

\paragraph{Findings.} On the validation targets V5.4 is worse than V5.3-A: 265 versus 280 successes. Paired by episode, V5.4 fixes 13 failures but breaks 28 successes (252 shared successes, 7 shared failures). An exact McNemar test on these 41 discordant episodes gives $p=0.028$, so the drop is unlikely to be chance. On the random-target, random-spawn episodes the two models are equivalent (76 versus 75, with almost entirely overlapping intervals), even though V5.4 was trained on exactly this distribution. The teacher-forced validation MSE agrees with the closed-loop result: it is 2.5$\times$ (HOLD) and 1.6$\times$ (REACH) higher for V5.4 (Fig.~\ref{dag:fig:v54train}).

\paragraph{Interpretation.} Each rollout contributes at most 3000 rows per phase, so every source is weighted by its number of episodes. In $\dagD_3$ the 1400 random-target rollouts outnumber the 161 rollouts on the deployment targets by almost 9:1, which shifts the arm fit away from the validation distribution; Eq.~\eqref{dag:obj} is only minimised on average over the aggregated mixture. Round 3 also does not address the base: the wheels were again driven by the expert, so the remaining step-cap failures, many of which are navigation failures (Sec.~\ref{dag:task}), are unaffected. \textbf{V5.3-A therefore remains the reported checkpoint}; the 300 test episodes are still untouched.
